%% kapa: exam=1 round=36 year=2025 period=2 form=A subject=accounting subject_ko=회계학 seq=32 number=72
%% source: Q-Net 감정평가사 기출문제 (cst00309) artlSeq=5231525 file=raw/1008/5231525/2025년 제36회 감정평가사 1차시험 2교시 문제지 원본.pdf
%% conversion: kordoc 4.24.51 -> tools/gen1.py (auto); review=proofread-sonnet; flags=none
\begin{question}{72}
(주)감평은 20×1년 초에 설립되었으며, 정상개별원가계산을 적용하고 있다. 다음은 20×1년 말 배부차이를 조정하기 전의 제조간접원가 계정과 기말 재고자산 및 매출원가에 관한 자료이다.
\begin{tcolorbox}[enhanced, colback=white, colframe=black, boxrule=0.4pt, arc=0pt, left=4pt, right=4pt, top=4pt, bottom=4pt]
- 제조간접원가
\begin{center}\small
\begin{tabular}{r|l}
\multicolumn{2}{c}{제조간접원가} \\ \hline
370,000\quad & \quad 310,000
\end{tabular}
\end{center}
- 기말 재고자산 및 매출원가
\begin{center}\small
\begin{tabular}{cccc}
\uline{\ 원재료\ } & \uline{\ 재공품\ } & \uline{\ 제품\ } & \uline{\ 매출원가\ } \\
￦500,000 & ￦200,000 & ￦300,000 & ￦1,500,000
\end{tabular}
\end{center}
\end{tcolorbox}
(주)감평은 제조간접원가 배부차이를 총원가비례배분법으로 조정하고 있다. 제조간접원가 배부차이를 조정한 후, (주)감평의 20×1년 매출원가는?
\begin{choices}
  \item ￦1,440,000
  \item ￦1,455,000
  \item ￦1,464,000
  \item ￦1,536,000
  \item ￦1,545,000
\end{choices}
\end{question}
\answer{5}
